\documentclass[12pt,a4paper]{report}

\usepackage[T1]{fontenc}
\usepackage[utf8]{inputenc}
\usepackage[french]{babel}
\usepackage{csquotes}
\usepackage{geometry}
\geometry{margin=2.5cm}
\usepackage{setspace}
\onehalfspacing
\usepackage{graphicx}
\usepackage{booktabs}
\usepackage{tabularx}
\usepackage{longtable}
\usepackage{array}
\usepackage{enumitem}
\usepackage{xcolor}
\usepackage{listings}
\usepackage{fancyhdr}
\usepackage{titlesec}
\usepackage{tocloft}
\usepackage{float}
\usepackage{hyperref}
\usepackage[nameinlink,noabbrev,french]{cleveref}
\usepackage[backend=biber,style=ieee,maxnames=6,sorting=none]{biblatex}
\addbibresource{biblio.bib}

\definecolor{ensiNavy}{HTML}{0B3D5C}

\hypersetup{
  colorlinks=true,
  linkcolor=ensiNavy,
  citecolor=ensiNavy,
  urlcolor=ensiNavy,
  pdftitle={LogiFlow TMS --- Rapport de projet tutoré},
  pdfauthor={ZOUAINE Oussama, STITOU Ibrahim, ES-SEMLLALI Oussama, RAHALI Hamza --- ENSI Tanger}
}

\lstset{
  basicstyle=\ttfamily\small,
  breaklines=true,
  frame=single,
  columns=fullflexible,
  keepspaces=true,
  showstringspaces=false
}

\newcommand{\repo}[1]{\texttt{#1}}
\newcommand{\blank}[1]{\underline{\hspace{#1}}}
\newcommand{\blankline}{\par\noindent\blank{12cm}\par}
% Figure flottante classique (diagrammes hauts)
\newcommand{\figmmd}[2]{%
  \begin{figure}[htbp]
    \centering
    \IfFileExists{figures/out/#1.png}{%
      \includegraphics[width=\linewidth,height=0.68\textheight,keepaspectratio]{figures/out/#1.png}%
    }{%
      \fbox{\parbox{0.9\linewidth}{\centering\textit{Diagramme #1 --- exécuter \texttt{scripts/render-mermaid.sh}}}}%
    }
    \caption{#2}
    \label{fig:#1}
  \end{figure}%
}
% Figure ancrée (diagrammes plats : pas de page quasi vide)
\newcommand{\figmmdhere}[2]{%
  \begin{figure}[H]
    \centering
    \vspace{-0.4em}
    \IfFileExists{figures/out/#1.png}{%
      \includegraphics[width=\linewidth,keepaspectratio]{figures/out/#1.png}%
    }{%
      \fbox{\parbox{0.9\linewidth}{\centering\textit{Diagramme #1 --- exécuter \texttt{scripts/render-mermaid.sh}}}}%
    }
    \caption{#2}
    \label{fig:#1}
    \vspace{-0.6em}
  \end{figure}%
}
% Capture ancrée [H] : évite les sections vides (figures qui flottent loin du titre)
\newcommand{\figshot}[2]{%
  \begin{figure}[H]
    \centering
    \vspace{-0.2em}
    \IfFileExists{figures/screenshots/#1.png}{%
      \includegraphics[width=0.88\linewidth,height=0.52\textheight,keepaspectratio]{figures/screenshots/#1.png}%
    }{%
      \fbox{\parbox[c][0.22\textheight][c]{0.9\linewidth}{%
        \centering
        \textbf{Capture à ajouter}\\[0.6em]
        \texttt{figures/screenshots/#1.png}\\[0.8em]
        \textit{#2}%
      }}%
    }
    \caption{#2}
    \label{fig:shot-#1}
    \vspace{-0.4em}
  \end{figure}%
}

\pagestyle{fancy}
\fancyhf{}
\fancyhead[L]{\small ENSI Tanger --- Projet tutoré}
\fancyhead[R]{\small LogiFlow TMS}
\fancyfoot[C]{\thepage}
\renewcommand{\headrulewidth}{0.4pt}

\begin{document}

%======================================================================
% PAGE DE GARDE — composition claire, hiérarchie typographique
%======================================================================
\begin{titlepage}
\thispagestyle{empty}
\setlength{\parindent}{0pt}

% --- Bandeau école ----------------------------------------------------
\begin{center}
  \IfFileExists{figures/logo-ensi.png}{%
    \includegraphics[height=2.55cm]{figures/logo-ensi.png}\\[0.55em]
  }{%
    {\color{ensiNavy}\large\bfseries École des Nouvelles Sciences \& Ingénierie}\\[0.35em]
  }
  {\large Tanger}
\end{center}

\vspace{1.1em}
{\color{ensiNavy}\hrule height 0.6pt}
\vspace{0.85em}

\begin{center}
  {\normalsize\bfseries\color{ensiNavy} RAPPORT DE PROJET TUTORÉ}\\[0.55em]
  {\small Filière Génie informatique --- temps aménagé}\\[0.2em]
  {\small Année universitaire 2025--2026}
\end{center}

\vspace{0.7em}
{\color{ensiNavy}\hrule height 0.35pt}

\vfill

% --- Bloc titre (signal principal) ------------------------------------
\begin{center}
  {\Huge\bfseries\color{ensiNavy} LogiFlow}\\[0.45em]
  {\Large\bfseries Transport Management System}\\[1em]
  \begin{minipage}{0.82\textwidth}
    \centering
    \large\color{ensiNavy!75}%
    Conception, architecture, intelligence artificielle\\[0.15em]
    et infrastructure d'un système de gestion\\[0.15em]
    du transport routier de marchandises
  \end{minipage}
\end{center}

\vfill

% --- Équipe -----------------------------------------------------------
{\color{ensiNavy}\hrule height 0.35pt}
\vspace{0.9em}

\begin{center}
  {\small\bfseries\color{ensiNavy} RÉALISÉ PAR}\\[0.75em]
  \begin{tabular}{@{}c@{\hspace{2.2cm}}c@{}}
    STITOU Ibrahim & ZOUAINE Oussama \\[0.55em]
    RAHALI Hamza & ES-SEMLLALI Oussama \\
  \end{tabular}
\end{center}

\vspace{1.15em}

\begin{center}
  {\small\bfseries\color{ensiNavy} ENCADRANT}\\[0.35em]
  {\large Jaber EL BOUHDIDI}
\end{center}

\vspace{0.6em}
{\color{ensiNavy}\hrule height 0.6pt}

\end{titlepage}
\clearpage
\setcounter{page}{1}


% Front matter — équipe LogiFlow

\chapter*{Dédicaces}
\addcontentsline{toc}{chapter}{Dédicaces}

Nous dédions ce travail à nos parents, dont le soutien, les sacrifices et la confiance ont rendu possible notre parcours à l'ENSI Tanger. Leur présence, dans les moments de doute comme dans les moments de réussite, a été le fondement de notre progression.

À nos familles et à nos proches, pour leur patience, leurs encouragements et leur bienveillance tout au long de cette formation.

À nos amis, compagnons de route, pour les échanges, l'entraide et la légèreté qu'ils ont su apporter durant les périodes les plus exigeantes du projet.

Enfin, nous dédions ce rapport à toute l'équipe LogiFlow : ce document est le fruit d'un travail collectif, de décisions partagées et d'un engagement commun pour concevoir un système utile, structuré et professionnel.

\vspace{1.2em}
\begin{flushright}
\textit{L'équipe LogiFlow}
\end{flushright}

\clearpage
\chapter*{Remerciements}
\addcontentsline{toc}{chapter}{Remerciements}

Nous tenons à remercier Allah, Le Tout-Puissant, de nous avoir accordé la force et la persévérance nécessaires pour mener à bien ce projet tutoré.

Nous exprimons notre profonde gratitude à notre encadrant, \textbf{Jaber EL BOUHDIDI}, pour son suivi, ses conseils et sa disponibilité tout au long du projet.

Nos remerciements vont à M.~\textbf{CHAKER Jihad}, responsable de la filière Génie informatique, ainsi qu'à l'ensemble des enseignants de l'École des Nouvelles Sciences et Ingénierie de Tanger pour la qualité de la formation reçue.

Enfin, nous remercions les membres du jury d'avoir accepté d'évaluer ce travail et de nous consacrer leur temps et leur expertise.

\clearpage
\chapter*{Résumé}
\addcontentsline{toc}{chapter}{Résumé}

Ce rapport présente le projet tutoré \textbf{LogiFlow}, un \emph{Transport Management System} (TMS) pour le transport routier de marchandises, réalisé en équipe de quatre étudiants ingénieurs à l'ENSI Tanger (filière Génie informatique, année 2025--2026).

L'objectif est de concevoir et de développer une plateforme web unique, filtrée par rôle métier (exploitant, responsable d'exploitation, atelier, commercial, chauffeur, administrateur), couvrant le cycle commercial $\rightarrow$ dossier de transport $\rightarrow$ voyage $\rightarrow$ suivi, carburant et maintenance.

Le backend est un monolithe modulaire \textbf{Spring Boot} / \textbf{Spring Modulith}, structuré en architecture hexagonale par module métier, avec PostgreSQL. Le frontend est une SPA \textbf{Angular} (signals, Tailwind). L'intelligence artificielle (copilote conversationnel SSE, planification assistée, calcul d'itinéraire, maintenance prédictive) est isolée dans un service \textbf{Flask} (LLM cloud compatible OpenAI, OSRM, pgvector), jamais exposé au navigateur : Spring agit comme BFF. L'infrastructure cloud AWS (une EC2, Caddy, Docker Compose) est provisionnée par \textbf{Terraform} et configurée par \textbf{Ansible}, sans SSH exposé (AWS SSM).

Ce travail a permis de mobiliser des compétences en conception métier, architecture logicielle, développement full-stack, intégration d'agents IA et infrastructure as code.

\vspace{1em}
\noindent\textbf{Mots-clés :} Projet tutoré, TMS, Transport, Spring Modulith, Architecture hexagonale, Angular, Keycloak, Flask, Groq, OSRM, Terraform, Ansible, AWS, Intelligence artificielle.

\clearpage
\chapter*{Abstract}
\addcontentsline{toc}{chapter}{Abstract}

This report presents the supervised team project \textbf{LogiFlow}, a road-freight \emph{Transport Management System} (TMS) developed by four engineering students at ENSI Tangier (Computer Engineering, academic year 2025--2026).

The goal is a single role-filtered web application covering the commercial order $\rightarrow$ transport dossier $\rightarrow$ trip $\rightarrow$ tracking, fuel and maintenance lifecycle.

The backend is a modular \textbf{Spring Boot} / \textbf{Spring Modulith} monolith with hexagonal modules and PostgreSQL. The frontend is an \textbf{Angular} SPA. AI features (SSE copilot, assisted planning, routing, predictive maintenance) run in an isolated \textbf{Flask} service with a cloud OpenAI-compatible LLM and \textbf{OSRM}; the browser never talks to Flask. AWS runs a single EC2 Docker stack behind Caddy, managed with \textbf{Terraform} and \textbf{Ansible} via SSM (no public SSH).

\vspace{1em}
\noindent\textbf{Keywords:} Supervised project, TMS, freight, Spring Modulith, hexagonal architecture, Angular, Keycloak, Flask, Groq, OSRM, Terraform, Ansible, AWS, Artificial Intelligence.

\clearpage
\chapter*{Avant-propos}
\addcontentsline{toc}{chapter}{Avant-propos}

\begin{description}[style=nextline,leftmargin=3.2cm,font=\normalfont\bfseries]
  \item[Établissement] École des Nouvelles Sciences \& Ingénierie (ENSI), Tanger
  \item[Filière] Génie informatique --- temps aménagé
  \item[Cadre] Projet tutoré (travail d'équipe)
  \item[Intitulé] LogiFlow --- Transport Management System
  \item[Encadrant] Jaber EL BOUHDIDI
\end{description}

\vspace{1.2em}
\noindent\textbf{Composition de l'équipe et responsabilités} :

\begin{table}[ht]
\centering
\small
\begin{tabular}{@{}clp{3.2cm}p{5.8cm}@{}}
\toprule
\# & Nom et prénom & Rôle(s) & Responsabilités principales \\
\midrule
1 & STITOU Ibrahim
  & Chef de projet ;\newline Responsable Backend
  & Coordination de l'équipe, suivi du planning, communication avec l'encadrant ;\newline développement API Spring Boot, gestion de la base de données \\[0.7em]
2 & ZOUAINE Oussama
  & Responsable Frontend ;\newline Responsable Infra / DevOps
  & Développement Angular, UI/UX, intégration API ;\newline Docker, CI/CD, déploiement AWS \\[0.7em]
3 & RAHALI Hamza
  & Responsable Sécurité
  & Authentification, gestion des accès, sécurité applicative \\[0.7em]
4 & ES-SEMLLALI Oussama
  & Responsable IA
  & Développement et intégration des fonctionnalités IA dans l'application \\
\bottomrule
\end{tabular}
\end{table}

\clearpage
\chapter*{Liste des sigles et acronymes}
\addcontentsline{toc}{chapter}{Liste des sigles et acronymes}

\begin{table}[ht]
\centering
\caption{Sigles et acronymes utilisés dans ce rapport.}
\label{tab:sigles}
\begin{tabular}{@{}lp{10cm}@{}}
\toprule
Acronyme & Désignation \\
\midrule
ADR & Architecture Decision Record \\
API & Application Programming Interface \\
BFF & Backend for Frontend \\
CI/CD & Continuous Integration / Continuous Delivery \\
DTO & Data Transfer Object \\
GHCR & GitHub Container Registry \\
IAM & Identity and Access Management \\
JWT & JSON Web Token \\
LLM & Large Language Model \\
MFA & Multi-Factor Authentication \\
OIDC & OpenID Connect \\
OSRM & Open Source Routing Machine \\
PKCE & Proof Key for Code Exchange \\
RBAC & Role-Based Access Control \\
SSE & Server-Sent Events \\
SSM & AWS Systems Manager Session Manager \\
TMS & Transport Management System \\
VPC & Virtual Private Cloud \\
\bottomrule
\end{tabular}
\end{table}


\clearpage
\listoffigures
\clearpage
\listoftables
\clearpage
\tableofcontents
\clearpage

\chapter{Introduction générale}
\label{chap:intro}

\section{Contexte}

Le transport routier de marchandises repose sur un enchaînement d'objets métier distincts : une demande commerciale, un dossier transportable et facturable, puis une exécution physique (véhicule, remorque, chauffeur, étapes de trajet). Sans système dédié, ces objets se mélangent dans des tableurs, les rôles (exploitation, atelier, commercial, chauffeur) n'ont pas la même vue, et les aides algorithmiques (itinéraire, groupage, assistant) restent collées au navigateur ou à une API LLM tierce.

LogiFlow est un \emph{Transport Management System} (TMS) conçu pour ce cycle, dans le cadre d'un \textbf{projet tutoré} à l'ENSI Tanger (équipe de quatre étudiants, filière Génie informatique). L'objectif est une base professionnelle : domaine explicite, API typée, discipline architecturale vérifiée en CI, agents IA isolés, et un déploiement cloud reproductible.

\section{Problématique}

Comment concevoir une plateforme TMS web unique, sécurisée et évolutive, qui :
\begin{enumerate}[leftmargin=*]
  \item sépare clairement commercial, transportable et exécution ;
  \item adapte l'interface aux six rôles métier sans multiplier les applications ;
  \item intègre l'IA (copilote, planification, itinéraire) sans exposer le service LLM au navigateur ;
  \item déploie sur le cloud avec une surface d'attaque maîtrisée ?
\end{enumerate}

\section{Objectifs}

\begin{enumerate}[leftmargin=*]
  \item Fournir une application web unique, filtrée par \textbf{rôle}.
  \item Encapsuler le métier dans un backend Java \textbf{modulaire} (Spring Modulith) et \textbf{hexagonal} au sein de chaque module \cite{evans2003,cockburn2005,springmodulith}.
  \item Isoler l'intelligence artificielle (LLM cloud compatible OpenAI, routing OSRM) derrière le backend : le navigateur ne parle jamais au service Flask.
  \item Authentifier via Keycloak (OIDC + PKCE, MFA côté IdP) et valider le JWT côté Spring.
  \item Décrire et provisionner l'infrastructure AWS par code (Terraform + Ansible), sans SSH exposé.
\end{enumerate}

\section{Organisation des dépôts}

\begin{table}[H]
\centering
\caption{Dépôts LogiFlow et rôle de chacun.}
\label{tab:repos}
\begin{tabular}{@{}ll@{}}
\toprule
Dépôt & Rôle \\
\midrule
\repo{logiflow-frontend} & SPA Angular (UI, auth OIDC, App Shell) \\
\repo{logiflow-backend} & API Spring Boot, PostgreSQL, façade IA \\
\repo{logiflow-ai-service} & Agents Flask + LLM cloud + OSRM \\
\repo{logiflow-infra} & Terraform (AWS) + Ansible (déploiement) \\
\bottomrule
\end{tabular}
\end{table}

Cette séparation permet de versionner l'infrastructure indépendamment du code : un \texttt{terraform apply} ne rebuild pas l'application ; la CI applicative publie une image Docker sur GHCR, Ansible ne fait que la \emph{pull}.

\section{Conduite du projet}

\subsection{Méthodologie}

Le projet suit une approche \textbf{agile} adaptée à une équipe de quatre : sprints courts, livraisons incrémentales (référentiels $\rightarrow$ cycle commande--dossier--voyage $\rightarrow$ IA $\rightarrow$ infra), revues régulières avec l'encadrant.

\subsection{Répartition des rôles}

\begin{table}[H]
\centering
\caption{Rôles et responsabilités de l'équipe LogiFlow.}
\label{tab:repartition}
\small
\begin{tabular}{@{}clp{3.4cm}p{5.6cm}@{}}
\toprule
\# & Membre & Rôle(s) & Responsabilités principales \\
\midrule
1 & STITOU Ibrahim
  & Chef de projet ;\newline Responsable Backend
  & Coordination de l'équipe, suivi du planning, communication avec l'encadrant ;\newline API Spring Boot, base de données \\[0.65em]
2 & ZOUAINE Oussama
  & Responsable Frontend ;\newline Responsable Infra / DevOps
  & Angular, UI/UX, intégration API ;\newline Docker, CI/CD, déploiement AWS \\[0.65em]
3 & RAHALI Hamza
  & Responsable Sécurité
  & Authentification, gestion des accès, sécurité applicative \\[0.65em]
4 & ES-SEMLLALI Oussama
  & Responsable IA
  & Développement et intégration des fonctionnalités IA \\
\bottomrule
\end{tabular}
\end{table}

\subsection{Planning}

Le diagramme de Gantt global croise \textbf{six volets} (cadre, conception, backend, frontend, IA, infra/qualité) avec des \textbf{jalons} de revue et de livraison. Les flux backend et frontend démarrent en parallèle après la revue de conception ; l'IA et l'infra s'appuient sur le socle API.

\figmmdhere{planning-gantt}{Diagramme de Gantt global du projet tutoré : sprints parallèles, dépendances et jalons.}

\begin{table}[H]
\centering
\caption{Jalons clés du planning.}
\label{tab:jalons}
\small
\begin{tabular}{@{}llp{7.2cm}@{}}
\toprule
Jalon & Cible & Livrable attendu \\
\midrule
Kick-off & Début projet & Périmètre, rôles équipe, outillage \\
Revue conception & Fin février & Glossaire, use cases, backlog priorisé \\
Démo intermédiaire & Fin mars & Cycle commande--voyage utilisable \\
Démo finale & Fin avril & IA + infra + parcours rôle complets \\
Dépôt rapport & Fin avril & PDF + sources + captures \\
\bottomrule
\end{tabular}
\end{table}

Les weekends sont exclus du Gantt pour coller au rythme scolaire.

\subsection{Plannings individuels}

Chaque membre dispose d'un Gantt aligné sur ses rôles, avec les mêmes jalons d'équipe (revue conception, démos, dépôt).

\figmmdhere{gantt-stitou}{Gantt individuel --- STITOU Ibrahim (chef de projet et responsable backend).}

\figmmdhere{gantt-zouaine}{Gantt individuel --- ZOUAINE Oussama (responsable frontend et infrastructure / DevOps).}

\figmmdhere{gantt-rahali}{Gantt individuel --- RAHALI Hamza (responsable sécurité).}

\figmmdhere{gantt-essemllali}{Gantt individuel --- ES-SEMLLALI Oussama (responsable intelligence artificielle).}

\section{Plan du rapport}

Le chapitre~\ref{chap:conception} pose le domaine et les exigences. Le chapitre~\ref{chap:architecture} donne la vue d'ensemble. Les chapitres~\ref{chap:frontend} à~\ref{chap:infra} détaillent frontend, backend, IA et infrastructure. Le chapitre~\ref{chap:securite} traite de la sécurité. Les annexes reprennent le glossaire métier et les décisions d'architecture (ADR).

\chapter{Conception métier}
\label{chap:conception}

\section{Problème}

Un TMS routier doit relier trois niveaux que les équipes mélangent souvent :

\begin{enumerate}[leftmargin=*]
  \item le \textbf{commercial} : une commande client, des lignes de marchandise, un statut ;
  \item le \textbf{transportable} : un dossier issu d'une commande confirmée, avec segments de chargement et de déchargement sur des sites géolocalisés ;
  \item l'\textbf{exécution} : un voyage (ressources + trajet), le suivi, le carburant, la maintenance.
\end{enumerate}

LogiFlow pose ces objets comme des agrégats distincts, chacun dans un module backend, plutôt qu'un enregistrement unique « mission ».

\figmmd{cycle-metier}{Cycle métier LogiFlow : du client au voyage, avec flotte, suivi, carburant et maintenance.}

\section{Acteurs et rôles}

Un \textbf{Utilisateur} est un compte applicatif du module \texttt{iam} (login, email, rôles, flag actif). Ce n'est ni un utilisateur Keycloak ni un enregistrement de mot de passe. Les rôles portés sont :

\begin{center}
\texttt{EXPLOITANT}, \texttt{RESPONSABLE\_EXPLOITATION}, \texttt{ATELIER},\\
\texttt{COMMERCIAL}, \texttt{CHAUFFEUR}, \texttt{ADMINISTRATEUR}.
\end{center}

Le JWT Keycloak expose ces rôles sous la forme \texttt{ROLE\_*} (ou claims \texttt{roles} / \texttt{realm\_access.roles}). Le \textbf{provisioning} d'un vrai compte n'est pas une inscription libre : seul un \texttt{ADMINISTRATEUR} crée un Utilisateur via \texttt{POST /api/v1/utilisateurs}.

\figmmd{use-cases}{Cas d'usage principaux par rôle applicatif.}

\section{Exigences fonctionnelles}

\begin{table}[ht]
\centering
\caption{Exigences fonctionnelles principales.}
\label{tab:ef}
\begin{tabular}{@{}cp{3.2cm}p{8cm}@{}}
\toprule
Id & Thème & Description \\
\midrule
EF-01 & Référentiels & Gérer sites, clients, marchandises, stations. \\
EF-02 & Commercial & Créer / confirmer des commandes et lignes. \\
EF-03 & Dossiers & Ouvrir un dossier depuis une commande confirmée. \\
EF-04 & Planning & Créer un voyage, affecter flotte et chauffeurs. \\
EF-05 & Suivi & Enregistrer les événements d'exécution. \\
EF-06 & Carburant & Saisir et valider des prises de carburant. \\
EF-07 & Maintenance & Plans, ordres de travail, scores de santé. \\
EF-08 & Documents & Attacher des fichiers aux entités. \\
EF-09 & IA & Copilote, propositions de planning, itinéraire. \\
EF-10 & IAM & Provisionner utilisateurs et rôles. \\
\bottomrule
\end{tabular}
\end{table}

\section{Exigences non fonctionnelles}

\begin{table}[ht]
\centering
\caption{Exigences non fonctionnelles.}
\label{tab:enf}
\begin{tabular}{@{}cp{3.2cm}p{8cm}@{}}
\toprule
Id & Thème & Description \\
\midrule
ENF-01 & Sécurité & JWT OIDC, MFA IdP, moindre privilège réseau. \\
ENF-02 & Modularité & Frontières Modulith + ArchUnit en CI. \\
ENF-03 & Isolation IA & Flask inaccessible depuis Internet. \\
ENF-04 & Erreurs & Problèmes HTTP normalisés (RFC~7807). \\
ENF-05 & IaC & Infra AWS reproductible (Terraform). \\
ENF-06 & Observabilité & Health checks, journal des interactions IA. \\
\bottomrule
\end{tabular}
\end{table}

\section{Cycle de vie du transport}

\begin{table}[ht]
\centering
\caption{Concepts métier centraux et module propriétaire.}
\label{tab:concepts}
\begin{tabular}{@{}llp{6.2cm}@{}}
\toprule
Concept & Module & Définition \\
\midrule
Commande & \texttt{order} & Demande commerciale d'un client. \\
Dossier de transport & \texttt{dossier} & Unité transportable et facturable. \\
Voyage & \texttt{planning} & Exécution physique : ressources pour 1..n dossiers. \\
Trajet & \texttt{planning} & Séquence ordonnée d'étapes. \\
Groupage & \texttt{planning} & Plusieurs dossiers sur un même voyage. \\
\bottomrule
\end{tabular}
\end{table}

Un \textbf{Site} est un point de chargement ou de déchargement géolocalisé (WGS-84). Une \textbf{Station} (carburant) n'est \emph{pas} un Site : décision d'architecture séparée (annexe~\ref{chap:adr}).

Un dossier porte des \textbf{segments} (au moins un chargement et un déchargement). Un voyage porte une \textbf{affectation} chauffeur (titulaire ou renfort) et éventuellement une remorque, utilisée pour le taux de remplissage.

\section{Référentiels associés}

\begin{itemize}[leftmargin=*]
  \item \textbf{Véhicule} (tracteur, porteur, fourgon) vs \textbf{Remorque} (semi-remorque, capacités).
  \item \textbf{Chauffeur} (référentiel RH / conformité) distinct du rôle applicatif \texttt{CHAUFFEUR}.
  \item \textbf{Marchandise} (catalogue, métadonnées ADR) référencée par id sur la ligne de commande.
  \item \textbf{Plan d'entretien} (préventif) vs \textbf{ordre de travail} (exécution) vs \textbf{score de santé} (statut dérivé côté serveur).
  \item \textbf{Document} : pièce jointe polymorphe (\texttt{typeEntite} + \texttt{entiteId}).
\end{itemize}

\section{Expérience utilisateur cible}

L'\textbf{App Shell} offre une sidebar de navigation filtrée par rôle, un pied de sidebar avec identité et déconnexion, et une barre utilitaire (fil d'Ariane, copilote IA, badge File du jour). Le \textbf{Tableau de bord} combine un aperçu (compteurs / statuts) et une \textbf{File du jour} (file de travail priorisée). Une palette de commandes (\texttt{Ctrl/Cmd+K}) accélère la navigation.

\section{Diagrammes de conception UML}

Cette section regroupe les vues UML utilisées pour figer le domaine avant le découpage technique (chapitres suivants).

\subsection{Diagramme de cas d'utilisation}

\figmmd{usecase-uml}{Diagramme de cas d'utilisation LogiFlow : acteurs (rôles) et cas d'usage du système.}

Le diagramme précédent (figure~\ref{fig:use-cases}) montrait les liens de façon compacte ; la figure~\ref{fig:usecase-uml} reprend une présentation plus proche d'un use-case UML classique (système encapsulé, acteurs autour, enchaînements métier commande~$\rightarrow$~dossier~$\rightarrow$~voyage).

\subsection{Diagramme de classes (domaine)}

\figmmd{class-metier}{Diagramme de classes métier (cœur TMS) : Client, Commande, Dossier, Voyage, flotte et satellites.}

Les associations suivent les règles du glossaire : une commande confirmée génère un dossier ; un voyage groupe un ou plusieurs dossiers ; chauffeur et véhicule sont liés au voyage via l'affectation (pas un agrégat unique « mission »).

Chaque lien est porté par un attribut de type \texttt{FK} sur la classe dépendante (\texttt{voyageId} et \texttt{chauffeurId} sur Affectation, \texttt{dossierIds} sur Voyage). Le losange plein marque un objet contenu dans son parent : ligne de commande, segment, affectation. \texttt{remorqueId} reste optionnel (cardinalité \texttt{0..1}).

\subsection{Diagramme de séquence métier}

\figmmd{seq-commande-voyage}{Séquence métier : création / confirmation de commande, ouverture de dossier, planification et démarrage de voyage.}

\subsection{Diagramme d'états}

Ces états sont volontairement synthétiques : chaque agrégat (commande, dossier, voyage) possède son propre enum de statut côté code ; le schéma ci-dessous montre le \emph{fil conducteur} perçu par l'exploitation.

\figmmdhere{etat-cycle}{Machine d'états du cycle commande~$\rightarrow$~dossier~$\rightarrow$~voyage jusqu'à livraison.}

\chapter{Architecture d'ensemble}
\label{chap:architecture}

\section{Principe directeur}

Trois règles structurent tout le système :

\begin{enumerate}[leftmargin=*]
  \item \textbf{Le navigateur ne parle qu'à Spring Boot.} Angular n'appelle jamais Flask. Le service IA n'a pas d'IP publique ; seuls les security groups internes l'autorisent depuis le backend.
  \item \textbf{Flask n'a pas la base TMS.} Tout le contexte métier est assemblé par Spring (modules via leur package \texttt{api}) puis envoyé dans le corps HTTP interne.
  \item \textbf{L'identité utilisateur s'arrête au backend.} Flask ne voit pas le JWT. Entre Spring et Flask, une clé d'API de service (\texttt{X-Internal-Api-Key}) suffit ; les outils du copilote rappellent Spring avec une clé de callback.
\end{enumerate}

\section{Vue conteneurs}

\figmmd{c4-containers}{Vue conteneurs LogiFlow : Angular, Keycloak, Spring Boot, PostgreSQL, Flask, LLM cloud et OSRM.}

\begin{table}[ht]
\centering
\caption{Technologies par conteneur.}
\label{tab:stack}
\begin{tabular}{@{}ll@{}}
\toprule
Conteneur & Stack \\
\midrule
Frontend & Angular 22, Tailwind 4, Zard UI, OIDC client \\
Backend & Java 21, Spring Boot, Modulith, JPA, Flyway \\
Base TMS & PostgreSQL (+ PostGIS) \\
Service IA & Python 3.13, Flask, LLM cloud (Groq\ldots), OSRM, pgvector \\
Identité & Keycloak (OIDC Authorization Code + PKCE) \\
Infra & AWS EC2, Terraform, Ansible, GHCR \\
\bottomrule
\end{tabular}
\end{table}

\section{Flux d'une requête}

\figmmd{flux-global}{Flux global d'une requête utilisateur, y compris le chemin IA interne.}

L'utilisateur atteint le frontend (nginx + Angular) en HTTPS. Le frontend appelle \texttt{/api/v1/...} sur Spring avec un JWT. Pour le copilote ou l'itinéraire, Spring authentifie, charge le contexte, appelle Flask sur le port interne 8000, journalise l'interaction, et renvoie un DTO déjà mis en forme. Les erreurs suivent RFC~7807 \cite{rfc7807}.

\section{Découpage « un processus, plusieurs frontières »}

Le backend est un \textbf{seul} JAR / une seule image Docker, mais \textbf{plusieurs} modules métier (chapitre~\ref{chap:backend}). L'IA est un processus séparé parce qu'elle a un cycle de vie différent (dépendances Python, timeouts LLM, streaming SSE). L'infra est un dépôt séparé parce qu'un changement de VPC ne doit pas exiger un commit applicatif.

\section{Environnements}

\begin{description}[style=nextline,leftmargin=1.2cm]
  \item[Local] Docker Compose (PostgreSQL, Keycloak, Flask). Frontend OIDC + backend profil \texttt{dev} : JWT Keycloak obligatoire, comme en production. Angular sur le port 4200 (proxy vers 8080).
  \item[Cloud (AWS)] Une EC2 \texttt{t3.large} dans \texttt{eu-west-3} (VPC public, Elastic IP, SG 80/443), stack Docker Compose derrière Caddy (TLS), Keycloak sur \texttt{auth.<domaine>}. Administration par AWS SSM, aucun port 22 ouvert \cite{awsssm}. Pas de NAT Gateway (choix coût).
\end{description}

\chapter{Frontend}
\label{chap:frontend}

\section{Stack et organisation}

Le frontend (\repo{logiflow-frontend}) est une SPA \textbf{Angular 22} : composants standalone, signals, zoneless, formulaires signal, \texttt{httpResource}. L'UI s'appuie sur \textbf{Tailwind CSS 4} et \textbf{Zard UI}. Cartes Leaflet et graphiques ECharts complètent le tableau de bord.

Les fonctionnalités sont organisées par dossiers métier (\texttt{commandes/}, \texttt{voyages/}, \texttt{maintenance/}, \ldots) avec routes lazy sous un \textbf{SignedInShell}.

\section{App Shell et navigation}

L'\textbf{App Shell} fixe la chrome : sidebar de navigation primaire, zone de contenu scrollable, barre utilitaire (fil d'Ariane, déclencheur Copilote IA, badge File du jour). L'identité et la déconnexion vivent dans le \textbf{pied de sidebar}, pas dans un menu avatar.

La navigation est dérivée d'une table de destinations filtrée par rôles (\texttt{destinationsForRoles}). Un utilisateur ne voit que les modules qu'il a le droit d'ouvrir.

\section{Contrôle d'accès (RBAC)}

Chaque route porte \texttt{data.roles}. Le garde \texttt{roleGuard} redirige vers
\texttt{/403} (Accès refusé) si le rôle de session ne correspond pas. Un utilisateur
authentifié mais non autorisé voit le 403 ; un utilisateur sans session est renvoyé
vers la connexion.

Le tableau~\ref{tab:rbac-ui} résume, pour quelques zones de l'UI, quels rôles métier
y ont typiquement accès (claims portés par le JWT Keycloak).

\begin{table}[H]
\centering
\caption{Zones de l'interface et rôles métier associés (exemples).}
\label{tab:rbac-ui}
\small
\begin{tabularx}{\linewidth}{@{}>{\bfseries\raggedright\arraybackslash}p{3.4cm} >{\raggedright\arraybackslash}X@{}}
\toprule
Zone UI & Rôles autorisés (libellé et claim JWT) \\
\midrule
Commandes / clients &
  Commercial (\texttt{COMMERCIAL}) ;\quad
  Administrateur (\texttt{ADMINISTRATEUR}) \\[0.65em]
Dossiers / voyages &
  Exploitant (\texttt{EXPLOITANT}) ;\quad
  Responsable d'exploitation (\texttt{RESPONSABLE\_EXPLOITATION}) \\[0.65em]
Maintenance &
  Atelier (\texttt{ATELIER}) ;\quad
  Responsable d'exploitation (\texttt{RESPONSABLE\_EXPLOITATION}) \\[0.65em]
Carburant (saisie) &
  Chauffeur (\texttt{CHAUFFEUR}) ;\quad
  plus validateurs côté exploitation \\[0.65em]
Utilisateurs &
  Administrateur (\texttt{ADMINISTRATEUR}) \\
\bottomrule
\end{tabularx}
\end{table}

\section{Authentification}

Le chemin nominal --- en local comme en cloud --- est \textbf{Keycloak} via
\texttt{angular-auth-oidc-client} (Authorization Code + \textbf{PKCE}). Le frontend est un
client public OIDC : après connexion IdP, il joint l'access token sur chaque appel
\texttt{/api/}. Spring agit uniquement comme \emph{resource server} (validation JWT, rôles) ;
il n'émet pas de jetons et ne stocke pas les mots de passe
(voir chapitre~\ref{chap:securite}).

En développement comme en production, le frontend s'authentifie auprès de Keycloak
(OIDC) et le backend (profil \texttt{dev} ou \texttt{prod}) exige un JWT valide sur
l'API métier.

\section{Écrans clés}

Les principaux écrans sont illustrés au chapitre~\ref{chap:captures} :

\begin{itemize}[leftmargin=*]
  \item \textbf{Tableau de bord} : Aperçu (compteurs / breakdowns de statut) puis File du jour.
  \item \textbf{Listes métier} : recherche, filtres, empty states avec CTA unique.
  \item \textbf{Fiches} : détail véhicule, voyage, dossier, avec documents et scores.
  \item \textbf{Nouveau voyage} : sélection de dossiers, calcul d'itinéraire IA, affectations.
  \item \textbf{Copilote} : panneau latéral conversationnel (pas une route dédiée en v1).
  \item \textbf{Palette de commandes} : modules, actions « Nouveau \ldots », références.
\end{itemize}

\chapter{Interfaces et captures d'écran}
\label{chap:captures}

Ce chapitre présente l'application Angular LogiFlow vue par un administrateur authentifié.
Les captures suivent le parcours d'usage : connexion et shell, tableau de bord, cycle
commande--dossier--voyage, flotte et exploitation, puis les deux usages principaux de l'IA.

\section{Connexion et App Shell}

L'entrée dans l'application passe par Keycloak (OIDC). Une fois authentifié, l'App Shell
fixe la navigation (sidebar filtrée par rôle) et les raccourcis globaux (Copilote, File du jour).

\figshot{01-connexion}{Écran de connexion : gateway vers Keycloak.}

\figshot{02-app-shell}{App Shell après connexion : modules, Copilote et File du jour visibles.}

\section{Tableau de bord}

Le tableau de bord synthétise les modules accessibles (compteurs), la file du jour
(exceptions prioritaires) et une répartition synthétique (ex.\ véhicules).

\figshot{03-tableau-de-bord}{Tableau de bord chargé : aperçu KPI, file du jour et synthèse.}

\section{Cycle commercial et transport}

Le cycle métier commence par les commandes, se poursuit par le dossier transportable,
puis le voyage (création assistée ou liste d'exécution avec carte OSRM).

\figshot{04-liste-commandes}{Liste des commandes (recherche et filtres).}

\figshot{05-fiche-dossier}{Fiche dossier de transport : segments, statut et marchandises.}

\figshot{06-nouveau-voyage}{Nouveau voyage : choix du mode planification assistée (IA).}

\figshot{07-liste-voyages}{Liste des voyages avec panneau d'itinéraire (OSRM).}

\section{Flotte, carburant et maintenance}

Les modules satellites couvrent le parc roulant, les prises de carburant et l'atelier.

\figshot{08-fiche-vehicule}{Fiche véhicule (ex.\ ME-001-CE) : identification, marque, poids et statut Disponible.}

\figshot{09-carburant}{Module carburant : prises et stations.}

\figshot{10-maintenance}{Module maintenance : ordres de travail et plans d'entretien.}

\section{Intelligence artificielle}

Deux écrans illustrent l'IA côté utilisateur : le copilote conversationnel (façade Spring)
et la planification assistée qui propose des voyages comparables.

\figshot{11-copilote}{Copilote en ligne : question métier et réponse structurée (voyages en cours).}

\figshot{12-planification-ia}{Planification assistée : propositions comparées et bouton de choix.}

\chapter{Backend}
\label{chap:backend}

\section{Monolithe modulaire}

Le backend (\repo{logiflow-backend}) est un monolithe \textbf{Spring Modulith} : un seul processus déployable, plusieurs modules métier aux frontières vérifiées en CI (ArchUnit + tests Modulith). Les modules ne se parlent que via leur package \texttt{api} ; \texttt{shared} est le seul module ouvert.

\figmmd{modules-backend}{Graphe de dépendances entre modules Spring Modulith (flèches vers \texttt{shared} omises).}

\begin{table}[ht]
\centering
\caption{Modules métier et responsabilité.}
\label{tab:modules}
\begin{tabular}{@{}lp{8.2cm}@{}}
\toprule
Module & Responsabilité \\
\midrule
\texttt{shared} & VOs, erreurs RFC~7807, audit, helpers sécurité \\
\texttt{iam} & Utilisateurs et rôles applicatifs \\
\texttt{referential} & Sites, clients, marchandises \\
\texttt{fleet} & Véhicules et remorques \\
\texttt{driver} & Chauffeurs, conformité, disponibilités \\
\texttt{order} & Commandes et lignes \\
\texttt{dossier} & Dossiers de transport, segments \\
\texttt{planning} & Voyages, trajets, affectations, conformité \\
\texttt{tracking} & Événements d'exécution \\
\texttt{maintenance} & OT, plans, sinistres, scores de santé \\
\texttt{carburant} & Prises, stations \\
\texttt{document} & Pièces jointes polymorphes \\
\texttt{ai} & Façade BFF vers Flask, journal des interactions \\
\texttt{config} & Security, CORS, async, scheduling \\
\bottomrule
\end{tabular}
\end{table}

\section{Architecture hexagonale par module}

Chaque module suit le même découpage : \texttt{domain} (agrégats, règles), \texttt{application} (cas d'usage), \texttt{infrastructure} (JPA, REST, adapters), \texttt{api} (contrat sortant vers les autres modules).

\figmmd{hexagonal}{Architecture hexagonale typique d'un module LogiFlow.}

\section{Modèle de données logique}

Les schémas SQL sont séparés par module (\texttt{commande}, \texttt{dossier}, \texttt{planning}, \ldots). Les liens inter-modules sont des UUID / JSON, pas des \texttt{@ManyToOne} JPA transverses.

\figmmd{er-logique}{Schéma entité-relation logique du cœur TMS et des satellites flotte / carburant / maintenance.}

\section{API et erreurs}

L'API publique est versionnée sous \texttt{/api/v1/**}. Les problèmes métier et techniques sont renvoyés au format \textbf{RFC~7807} (\texttt{application/problem+json}) \cite{rfc7807}. Les migrations de schéma passent par \textbf{Flyway}.

\section{Événements inter-modules}

Les publications utilisent \texttt{ApplicationEventPublisher} et \texttt{@ApplicationModuleListener} (asynchrone, après commit). Exemples : confirmation de commande $\rightarrow$ ouverture de dossier ; immobilisation atelier $\rightarrow$ indisponibilité flotte pour le planning.

\chapter{Intelligence artificielle}
\label{chap:ia}

\section{Rôle du service}

Les agents IA vivent dans le dépôt \repo{logiflow-ai-service} : une application \textbf{Flask~3}
(Python~3.13, Gunicorn), appelée \emph{uniquement} par Spring Boot. Le navigateur n'appelle jamais
ce service : Spring agit comme BFF et relaie notamment le flux SSE du copilote.

Quatre capacités sont exposées sous le préfixe \texttt{/internal/ai/v1} :

\begin{itemize}[leftmargin=*]
  \item \textbf{Copilote} --- conversations persistées, tool-calling vers Spring, RAG optionnel (pgvector), SSE ;
  \item \textbf{Planification} --- solveur déterministe de voyages + rédaction LLM ;
  \item \textbf{Maintenance prédictive} --- analyse déterministe (scores, anomalies) + synthèse LLM ;
  \item \textbf{Itinéraire} --- distance / durée / segments via \textbf{OSRM} uniquement \cite{osrm}.
\end{itemize}

\figmmd{ia-architecture}{Architecture du service IA : Spring BFF, agents Flask, LLM cloud, OSRM et base \texttt{logiflow\_ai}.}

\section{Organisation du code}

\begin{table}[ht]
\centering
\caption{Organisation du dépôt \repo{logiflow-ai-service}.}
\label{tab:ia-layout}
\small
\begin{tabular}{@{}lp{9.2cm}@{}}
\toprule
Zone & Contenu \\
\midrule
\texttt{app.py} / \texttt{wsgi.py} & Factory Flask, \texttt{GET /health}, point d'entrée Gunicorn \\
\texttt{api/v1/} & Blueprints protégés (\texttt{copilot}, \texttt{planification}, \texttt{maintenance}, \texttt{itinerary}) \\
\texttt{agents/} & Orchestrateurs et solveurs métier par agent \\
\texttt{infrastructure/} & \texttt{LlmClient}, \texttt{OsrmClient}, \texttt{BackendClient}, SQLAlchemy, repos \\
\texttt{security.py} & Vérification \texttt{X-Internal-Api-Key} (\texttt{hmac.compare\_digest}) \\
\texttt{migrations/} & Alembic (schéma \texttt{copilote}, extension \texttt{vector}) \\
\bottomrule
\end{tabular}
\end{table}

\section{Principes de conception}

\begin{enumerate}[leftmargin=*]
  \item \textbf{Pas d'accès à la base TMS.} Planification, maintenance et itinéraire reçoivent tout le contexte dans le corps HTTP. Le copilote lit le métier uniquement via des \textbf{outils} Spring.
  \item \textbf{Base propre} \texttt{logiflow\_ai} (PostgreSQL + \textbf{pgvector}, schéma \texttt{copilote}) : conversations, messages, appels d'outils, feedback, fragments de connaissance.
  \item \textbf{Pas de JWT / Keycloak} dans Flask. Spring pose \texttt{X-Utilisateur-Id} et \texttt{X-Copilote-Contexte} pour les rappels d'outils.
  \item \textbf{LLM cloud compatible OpenAI} (Groq par défaut via \texttt{LLM\_BASE\_URL} / \texttt{LLM\_API\_KEY} / \texttt{LLM\_MODEL}). Un LLM local (Ollama) a été écarté : trop lent sur CPU pour le streaming \cite{groq,ollama}.
  \item \textbf{Calcul déterministe, LLM rédacteur.} Chiffres, horaires et scores ne dépendent jamais du LLM ; un gabarit texte prend le relais si le fournisseur est indisponible.
\end{enumerate}

\section{Authentification Spring $\leftrightarrow$ Flask}

\begin{table}[ht]
\centering
\caption{Clés et en-têtes entre Spring et le service IA.}
\label{tab:ia-auth}
\begin{tabularx}{\linewidth}{@{}l X X@{}}
\toprule
Sens & En-tête / secret & Variable \\
\midrule
Spring $\rightarrow$ Flask & \texttt{X-Internal-Api-Key} &
  Flask \texttt{INTERNAL\_API\_KEY} $=$ Spring \texttt{AI\_SERVICE\_API\_KEY} \\
Flask $\rightarrow$ Spring (outils) & \texttt{X-Internal-Api-Key} (clé distincte) &
  Flask \texttt{BACKEND\_CALLBACK\_API\_KEY} $=$ Spring \texttt{AI\_SERVICE\_CALLBACK\_API\_KEY} \\
Contexte outils & \texttt{X-Copilote-Contexte} & Jeton opaque fourni par Spring \\
Identité UI & \texttt{X-Utilisateur-Id} & Posé après validation JWT \\
\bottomrule
\end{tabularx}
\end{table}

La route publique \texttt{GET /health} (hors préfixe v1) expose l'état du processus, du LLM
(\texttt{UP}, \texttt{CLE\_ABSENTE}, \texttt{DOWN}\ldots) et de la base, sans clé --- utilisée par le
front (badge « En ligne ») et les healthchecks de déploiement.

\section{Agents et contrats}

\begin{table}[ht]
\centering
\caption{Agents IA : routes Flask internes et exposition Spring (navigateur).}
\label{tab:agents}
\renewcommand{\arraystretch}{1.35}
\begin{tabularx}{\linewidth}{@{}>{\bfseries}l X@{}}
\toprule
Agent & Endpoints \\
\midrule
Copilote &
  \textbf{Flask :} \texttt{/internal/ai/v1/copilot/conversations/**}\newline
  \quad (messages en \texttt{text/event-stream}, SSE)\newline
  \textbf{Spring :} \texttt{/api/v1/ia/copilote/**} \\
Planification &
  \textbf{Flask :} \texttt{POST .../planification/proposer}\newline
  \textbf{Spring :} \texttt{POST .../ia/planification/propositions} \\
Itinéraire &
  \textbf{Flask :} \texttt{POST .../itinerary/calculer}\newline
  \textbf{Spring :} \texttt{POST .../ia/itineraires/calcul} \\
Maintenance &
  \textbf{Flask :} \texttt{POST .../maintenance/recommander}\newline
  \textbf{Spring :} \texttt{POST .../ia/maintenance/analyse} \\
\bottomrule
\end{tabularx}
\end{table}

\begin{table}[ht]
\centering
\caption{Répartition déterministe / LLM par agent.}
\label{tab:ia-det-llm}
\small
\begin{tabular}{@{}lll@{}}
\toprule
Agent & Déterministe & LLM \\
\midrule
Copilote & Persistance SQL, exécution d'outils & Boucle tool-calling + streaming \\
Planification & Groupage, matrice OSRM/Haversine, horaires & Comparaison et justifications \\
Maintenance & Scores, échéances, anomalies & Synthèse textuelle \\
Itinéraire & OSRM \texttt{/route} uniquement & --- \\
\bottomrule
\end{tabular}
\end{table}

\section{Séquence copilote}

Spring ouvre un flux SSE ; Flask charge l'historique, récupère le catalogue d'outils
(\texttt{GET /internal/copilote/outils}), boucle jusqu'à \texttt{COPILOTE\_MAX\_ITERATIONS\_OUTILS}
(défaut~4), exécute les outils via Spring (\texttt{POST /internal/copilote/outils/\{nom\}}),
puis streame les jetons. Événements SSE : \texttt{meta}, \texttt{attente}, \texttt{outil},
\texttt{token}, \texttt{sources}, \texttt{titre}, \texttt{fin}, \texttt{erreur}.

Outils métier typiques (Spring) : recherche dossiers / voyages / commandes / clients / chauffeurs,
proposition de voyages, consultation maintenance, consommation carburant. Outil local Flask
(si embeddings configurés) : \texttt{rechercher\_base\_connaissance}.

\figmmd{seq-copilote}{Séquence du copilote : Angular $\rightarrow$ Spring $\rightarrow$ Flask $\rightarrow$ outils Spring $\rightarrow$ LLM cloud (SSE).}

\section{Planification assistée}

Le solveur groupe les dossiers candidats, calcule une matrice de distances OSRM
(repli Haversine~$\times$~1,3), pose créneaux / pauses / équipage, puis le LLM rédige
justifications et comparaison (\texttt{source\_redaction}). Spring \textbf{réévalue la conformité}
avant restitution à l'UI ; la création du voyage reste une mutation TMS explicite.

\figmmd{seq-planification}{Séquence de planification assistée avec contrôle de conformité côté Spring.}

\section{Itinéraire et maintenance}

\begin{itemize}[leftmargin=*]
  \item \textbf{Itinéraire} : client OSRM HTTP ; réponse \textbf{503} si OSRM indisponible (aucun trajet inventé).
  \item \textbf{Maintenance} : score de risque, échéances de plans OT, anomalies (documents, sinistralité, surconsommation) calculés en déterministe ; le LLM synthétise ; gabarit si LLM down.
\end{itemize}

\section{Schéma de données copilote}

La base \texttt{logiflow\_ai} est distincte de la base TMS. Le schéma PostgreSQL \texttt{copilote}
est géré par Alembic (extension \texttt{vector}).

\figmmd{ia-schema}{Schéma relationnel du domaine copilote (conversations, messages, outils, RAG).}

\begin{itemize}[leftmargin=*]
  \item \texttt{conversation} --- fil de discussion par \texttt{utilisateur\_id} ;
  \item \texttt{message} --- rôles \texttt{user}/\texttt{assistant}, \texttt{sources} JSONB, tokens, modèle ;
  \item \texttt{appel\_outil} --- journal des tool-calls (nom, arguments, succès, durée) ;
  \item \texttt{feedback} --- note optionnelle sur un message ;
  \item \texttt{document\_connaissance} / \texttt{fragment\_connaissance} --- RAG avec colonne \texttt{embedding}.
\end{itemize}

\section{Stack et configuration}

\begin{table}[ht]
\centering
\caption{Stack du service \repo{logiflow-ai-service}.}
\label{tab:ia-stack}
\begin{tabular}{@{}ll@{}}
\toprule
Composant & Choix \\
\midrule
Runtime & Python 3.13, uv, Flask 3, Gunicorn (gthread, SSE) \\
Validation & Pydantic v2 (snake\_case Python / camelCase JSON) \\
Persistance & SQLAlchemy 2, Alembic, PostgreSQL + pgvector \\
LLM & API compatible OpenAI (Groq par défaut) \\
HTTP sortant & httpx (Spring, OSRM, LLM) \\
Routing & OSRM (\texttt{OSRM\_BASE\_URL}) \\
\bottomrule
\end{tabular}
\end{table}

Variables clés : \texttt{INTERNAL\_API\_KEY}, \texttt{DATABASE\_URL}, \texttt{LLM\_*},
\texttt{EMBED\_*} (optionnel), \texttt{BACKEND\_BASE\_URL}, \texttt{BACKEND\_CALLBACK\_API\_KEY},
\texttt{COPILOTE\_MAX\_ITERATIONS\_OUTILS}, \texttt{COPILOTE\_BATTEMENT\_S}, \texttt{OSRM\_*}.

\chapter{Infrastructure}
\label{chap:infra}

\section{Objectifs et choix structurants}

Le dépôt \repo{logiflow-infra} provisionne et déploie un environnement de démonstration sur AWS
(\texttt{eu-west-3}). Principes :

\begin{itemize}[leftmargin=*]
  \item infrastructure as code (\textbf{Terraform}) et configuration as code (\textbf{Ansible}) ;
  \item images applicatives construites en CI et poussées sur \textbf{GHCR} --- jamais buildées sur le serveur ;
  \item administration \textbf{sans SSH public} (AWS Systems Manager) \cite{awsssm} ;
  \item \textbf{une seule EC2} (\texttt{t3.large}) hébergeant toute la stack Docker Compose
        (ADR « instance unique ») --- choix volontaire pour contenir les coûts étudiants,
        plutôt que trois machines FE / BE / AI.
\end{itemize}

\section{Topologie AWS}

\begin{itemize}[leftmargin=*]
  \item VPC \texttt{10.20.0.0/16}, \textbf{un sous-réseau public}, Internet Gateway, Elastic IP ;
  \item \textbf{pas de NAT Gateway} (économie ~35~\$/mois) : la sortie Internet se fait depuis l'EC2 publique ;
  \item security group \texttt{logiflow-prod-serveur} : ingress \textbf{TCP 80 et 443} uniquement ;
        aucun port 22 ;
  \item reverse proxy \textbf{Caddy} (TLS Let's Encrypt, HTTP/2--3) devant frontend, backend et Keycloak ;
  \item conteneurs internes : Angular/nginx, Spring Boot, Flask IA, Keycloak, PostgreSQL
        (PostGIS + pgvector, trois bases : \texttt{logiflow}, \texttt{logiflow\_ai}, \texttt{keycloak}).
\end{itemize}

\figmmd{infra-aws}{Topologie AWS réelle : une EC2 publique, Caddy, stack Docker, SSM, GHCR et LLM cloud.}

\subsection{Routage Caddy (même origine)}

L'UI et l'API partagent \texttt{https://app.<domaine>} : pas de CORS navigateur.
\texttt{/api/*} est proxifié vers Spring avec \texttt{flush\_interval -1} pour le SSE copilote.
Keycloak est servi sur \texttt{https://auth.<domaine>}.

\section{Modules Terraform}

\begin{table}[ht]
\centering
\caption{Modules Terraform de \repo{logiflow-infra}.}
\label{tab:tf-modules}
\small
\begin{tabular}{@{}lp{8.8cm}@{}}
\toprule
Module & Ressources principales \\
\midrule
\texttt{bootstrap} & État S3 (verrou natif), OIDC GitHub, rôles IAM CI, zone Route~53 \\
\texttt{reseau} & VPC, subnet public, IGW, table de routes \\
\texttt{serveur} & EC2 Ubuntu 24.04, EIP, SG, instance profile SSM, alarme CloudWatch \\
\texttt{dns} & Enregistrements A \texttt{app}, \texttt{auth}, racine, \texttt{www} \\
\texttt{secrets} & SSM Parameter Store \texttt{/logiflow/prod/secrets/*} et \texttt{/config/*} \\
\texttt{sauvegardes} & Buckets S3 (rétention courte) + scripts de dump \\
\texttt{couts} & AWS Budgets, EventBridge Scheduler (arrêt auto 20~h) \\
\texttt{securite} & GuardDuty, Access Analyzer, CloudTrail, VPC Flow Logs, SNS \\
\bottomrule
\end{tabular}
\end{table}

\section{Ansible et stack runtime}

Playbook \texttt{site.yml} : rôles \texttt{base} $\rightarrow$ \texttt{docker} $\rightarrow$
\texttt{securite} $\rightarrow$ \texttt{logiflow} $\rightarrow$ \texttt{sauvegardes}.
Le rôle \texttt{logiflow} lit SSM, dépose la stack sous \texttt{/opt/logiflow}, authentifie GHCR
et exécute \texttt{bin/deployer.sh} (\texttt{docker compose pull} + healthchecks).

Images typiques : \texttt{ghcr.io/.../logiflow-\{frontend,backend,ai-service,keycloak,postgres\}}.
Mémoires plafonnées (~4,5~Go utiles sur 8~Go) pour cohabiter sur la \texttt{t3.large}.

\section{Secrets et IAM}

\begin{itemize}[leftmargin=*]
  \item Secrets générés par Terraform (\texttt{random\_password}) : DB, Keycloak, clés API internes AI ;
  \item stockage dans \textbf{SSM Parameter Store} (SecureString), jamais dans Git ;
  \item clé LLM injectée hors apply (\texttt{lifecycle.ignore\_changes}) ;
  \item CI GitHub via \textbf{OIDC} uniquement (rôles \texttt{logiflow-github-terraform} et
        \texttt{logiflow-github-deploiement}) --- pas de clés AWS longue durée.
\end{itemize}

\section{Pipeline de déploiement}

\figmmd{deploy}{Enchaînement Terraform (OIDC) $\rightarrow$ images GHCR $\rightarrow$ Ansible / SSM Run Command.}

\begin{enumerate}[leftmargin=*]
  \item Bootstrap ponctuel (bucket state, OIDC, Route~53).
  \item \texttt{terraform apply} sur \texttt{environments/prod} (manuel, environnement GitHub protégé).
  \item Build CI des dépôts applicatifs $\rightarrow$ tags GHCR.
  \item Workflow \texttt{deployer.yml} : démarre l'EC2 si arrêtée, \texttt{ssm:SendCommand},
        pull Compose, healthchecks.
\end{enumerate}

\section{CI infra}

\begin{itemize}[leftmargin=*]
  \item \texttt{qualite.yml} : \texttt{terraform fmt/validate}, tflint, ansible-lint, shellcheck, validation Compose/Caddy ;
  \item \texttt{terraform.yml} : plan commenté sur PR ; apply manuel ;
  \item \texttt{deployer.yml} : déploiement SSM ;
  \item \texttt{securite.yml} / \texttt{dast.yml} : Gitleaks, Trivy, CodeQL, ZAP.
\end{itemize}

\section{Coûts}

L'enveloppe cible est de l'ordre de 70~\$/mois. Des postes fixes subsistent même EC2 arrêtée
(EBS, EIP, Route~53, alarmes). Leviers principaux : \textbf{pas de NAT}, arrêt automatique à 20~h
(EventBridge), éventuel passage en \texttt{t3.medium}, rétention S3 courte. Le coût du fournisseur
LLM reste proportionnel à l'usage.

\chapter{Sécurité}
\label{chap:securite}

\section{Authentification et autorisation}

LogiFlow sépare clairement \textbf{émission} et \textbf{consommation} des jetons.
Keycloak est le fournisseur d'identité (IdP) OIDC ; Spring Boot est un
\emph{resource server} qui \textbf{valide} les JWT sans jamais stocker de mots de passe
ni émettre de jetons lui-même \cite{keycloak}.

\figmmd{seq-auth}{Séquence d'authentification : Angular $\leftrightarrow$ Keycloak $\leftrightarrow$ Spring Boot.}

\begin{itemize}[leftmargin=*]
  \item Angular : client OAuth2 / OIDC (Authorization Code + \textbf{PKCE}) via
        \texttt{angular-auth-oidc-client}.
  \item Spring : validation JWT (signature JWKS, issuer public, audience, rôles) ;
        clés lues via le réseau Docker interne vers Keycloak.
  \item MFA : entièrement gérée par l'IdP ; Spring peut exiger les claims \texttt{amr} / \texttt{acr}.
  \item RBAC applicatif : six rôles métier portés dans le JWT et contrôlés aussi côté UI.
  \item Deny-by-default sur l'API : seuls health / info / OpenAPI (désactivable) restent publics.
\end{itemize}

\figmmd{auth-responsabilites}{Répartition des responsabilités : Keycloak émet, Angular transporte, Spring autorise, Flask reste hors JWT utilisateur.}

\section{Pourquoi Keycloak (et pas « seulement du JWT »)}

\subsection{JWT seul ne suffit pas}

Un JWT n'est qu'un \textbf{format} de jeton signé. Il ne répond pas, à lui seul, aux questions
opérationnelles de l'identité :

\begin{itemize}[leftmargin=*]
  \item Qui crée les comptes, réinitialise les mots de passe, gère les sessions et la révocation ?
  \item Où vit le MFA (OTP, WebAuthn) et qui refuse d'émettre un jeton tant que le second facteur n'est pas validé ?
  \item Comment le frontend obtient-il le jeton sans embarquer un secret client (d'où PKCE) ?
\end{itemize}

Si Spring \emph{émettait} lui-même des JWT après un login login/mot de passe maison, le TMS
deviendrait un mini-IdP : stockage de credentials, politiques de mot de passe, récupération de
compte, MFA, audit --- hors du cœur métier transport. C'est l'anti-pattern écarté.

\subsection{Ce que Keycloak apporte}

\begin{enumerate}[leftmargin=*]
  \item \textbf{Standards} --- OIDC / OAuth2 prêts à l'emploi (Authorization Code + PKCE pour SPA).
  \item \textbf{IdP spécialisé} --- realm \texttt{logiflow}, clients, rôles, thèmes, import de
        configuration versionnée (image Docker dédiée).
  \item \textbf{MFA déléguée} --- aucun OTP en base TMS ; le backend peut seulement \emph{exiger}
        la preuve dans le jeton.
  \item \textbf{SSO possible} --- un même IdP peut servir d'autres applications plus tard.
  \item \textbf{Auto-hébergé et open source} --- adapté à un projet tutoré / démo cloud maîtrisée
        (même EC2, derrière Caddy sur \texttt{auth.<domaine>}), sans dépendance exclusive à un SaaS.
\end{enumerate}

\subsection{Alternatives examinées}

\begin{table}[H]
\centering
\caption{Alternatives d'identité et motif de non-retenue (ou usage limité).}
\label{tab:auth-alt}
\small
\begin{tabularx}{\linewidth}{@{}l X@{}}
\toprule
Option & Analyse \\
\midrule
Login + JWT maison (Spring) &
  Couplage fort ; Spring devient IdP ; MFA et reset à réécrire ; écarté. \\
Session serveur (cookie) &
  Moins adapté à une API stateless multi-clients ; scaling et SPA moins naturels. \\
Auth0 / Cognito / Firebase Auth &
  IdP matures en SaaS, mais coût, vendor lock-in et dépendance réseau externe ;
  moins de contrôle du realm pour une démo scolaire auto-hébergée. \\
Keycloak &
  \textbf{Retenu et utilisé} (local \texttt{dev} et production) :
  OIDC, MFA, auto-hébergeable, SPA + resource server Spring. \\
\bottomrule
\end{tabularx}
\end{table}

\figmmdhere{auth-choix}{Arbre de décision simplifié ayant conduit à Keycloak.}

\subsection{Conséquences pour LogiFlow}

\begin{itemize}[leftmargin=*]
  \item Le module \texttt{iam} du backend gère le \emph{profil applicatif} (rôles métier TMS),
        pas les secrets de connexion.
  \item Flask ne voit jamais le JWT utilisateur : après authentification Spring, seuls des
        en-têtes de service et un contexte copilote borné traversent la frontière IA.
  \item En cloud, l'\texttt{iss} du jeton est l'URL \textbf{publique} Keycloak ; la validation
        JWKS peut passer par le hostname Docker interne --- aucun aller-retour inutile par Internet
        pour les clés.
\end{itemize}

\section{Sécurité réseau et cloud}

\begin{itemize}[leftmargin=*]
  \item Surface d'attaque réduite : SG ouvert uniquement sur \textbf{80/443} ; reverse proxy Caddy ;
        backend, IA, Keycloak et Postgres \textbf{non exposés} hors réseau Docker.
  \item Pas de SSH ; administration via AWS SSM Session Manager / Run Command \cite{awsssm}.
  \item Secrets dans \textbf{SSM Parameter Store}, jamais dans Git ; instance profile en moindre privilège.
  \item EBS chiffré, IMDSv2, durcissement OS (rôle Ansible \texttt{securite} : CrowdSec, auditd\ldots).
  \item CI : tflint, scans (Gitleaks, Trivy, CodeQL) ; \texttt{terraform apply} et déploiement manuels
        (environnement GitHub \texttt{production}).
\end{itemize}

\section{Frontière Spring $\leftrightarrow$ Flask}

Flask ne reçoit pas le JWT utilisateur. L'appel entrant exige \texttt{X-Internal-Api-Key}.
Les rappels d'outils du copilote utilisent une clé de callback distincte et un contexte
\texttt{X-Copilote-Contexte}. Évolution possible vers mTLS ou OAuth2 client-credentials.

\section{Mesures côté code}

\begin{itemize}[leftmargin=*]
  \item Frontières Modulith + ArchUnit.
  \item Domaine sans dépendance framework.
  \item Erreurs HTTP normalisées (RFC~7807).
  \item Validation d'arguments sur les endpoints publics.
\end{itemize}

\section{Limites et axes d'amélioration}

L'architecture décrite est mise en œuvre dans les dépôts. Les axes suivants restent ouverts
pour une industrialisation ultérieure :

\begin{enumerate}[leftmargin=*]
  \item Renforcer Spring $\leftrightarrow$ Flask (mTLS ou OAuth2 client-credentials) au-delà des clés API.
  \item Approfondir les agents (géométrie de groupage, maintenance prédictive avancée).
  \item Étendre le monitoring et les alertes au-delà des healthchecks et budgets AWS.
  \item Maîtriser le coût LLM (quotas fournisseur) et le calendrier d'arrêt de l'EC2.
  \item Généraliser la MFA Keycloak sur l'ensemble des comptes cloud.
\end{enumerate}

Ces évolutions s'appuient sur les frontières déjà en place : BFF, modules Modulith, security groups,
SSM, Caddy et séparation des bases.

\chapter{Conclusion et perspectives}
\label{chap:conclusion}

\section{Bilan}

Ce projet tutoré a permis de concevoir LogiFlow, un TMS routier cohérent sur le plan
architectural : domaine explicite (commande $\rightarrow$ dossier $\rightarrow$ voyage),
monolithe modulaire hexagonal, SPA Angular filtrée par rôle, agents IA isolés derrière un
BFF Spring, et infrastructure AWS reproductible (une EC2, Terraform, Ansible, SSM) sans SSH exposé.

Les diagrammes de ce rapport synthétisent les décisions structurantes : isolation de Flask,
validation JWT Keycloak, graphe Modulith, topologie VPC / Caddy / Compose, et pipeline
GHCR $\rightarrow$ SSM.

\section{Apports pour l'équipe}

\begin{itemize}[leftmargin=*]
  \item Conception métier et glossaire partagé.
  \item Architecture logicielle (DDD léger, hexagonal, Modulith).
  \item Full-stack (Angular + Spring + PostgreSQL).
  \item Intégration d'agents IA (Flask, LLM cloud, OSRM, outils Spring) avec contraintes de sécurité.
  \item Infrastructure as code et culture DevSecOps (OIDC GitHub, SSM, scans CI).
\end{itemize}

\section{Difficultés rencontrées}

Plusieurs frictions ont structuré nos arbitrages. Brancher correctement Keycloak
(OIDC + PKCE) de bout en bout --- frontend, issuer public, JWKS interne, profils Spring
\texttt{dev}/prod --- a demandé une configuration stricte pour que le local reproduise le
cloud. Côté IA, la latence d'un LLM local a conduit à adopter un fournisseur cloud et à
séparer calcul déterministe / rédaction ; le streaming SSE a dû être préservé de bout en
bout (Caddy \texttt{flush\_interval}, Gunicorn gthread). Sur l'infrastructure, le budget
étudiant a motivé le choix d'une seule EC2 sans NAT, au prix d'une surface réseau différente
d'une architecture multi-tiers classique --- documentée dans les ADR. Enfin, la coordination
de quatre dépôts (front, back, IA, infra) a nécessité des contrats d'API stables et des
revues régulières avec l'encadrant.

\section{Perspectives}

\begin{enumerate}[leftmargin=*]
  \item Activer MFA Keycloak de bout en bout sur tous les comptes cloud.
  \item Renforcer Spring $\leftrightarrow$ Flask (mTLS ou OAuth2 client-credentials).
  \item Enrichir le copilote (outils, RAG) et la planification (contraintes métier avancées).
  \item Extraire un module (\texttt{planning} ou \texttt{tracking}) en service séparé
        \emph{seulement} si la charge le justifie --- les frontières \texttt{api} sont déjà prêtes.
  \item Industrialiser le monitoring (métriques, alertes) au-delà des healthchecks et budgets AWS.
\end{enumerate}

LogiFlow constitue ainsi une base solide pour poursuivre le développement et une
industrialisation progressive.


\appendix
\chapter{Glossaire métier}
\label{chap:glossaire}

\begin{description}[style=nextline,leftmargin=1.2cm,font=\normalfont\bfseries]
  \item[Utilisateur] Compte applicatif IAM (login, email, rôles, actif). Pas un utilisateur Keycloak.
  \item[Rôle] Une des valeurs \texttt{EXPLOITANT}, \texttt{RESPONSABLE\_EXPLOITATION}, \texttt{ATELIER}, \texttt{COMMERCIAL}, \texttt{CHAUFFEUR}, \texttt{ADMINISTRATEUR}.
  \item[Site] Point de chargement / déchargement géolocalisé.
  \item[Station] Point de distribution de carburant (distinct d'un Site).
  \item[Commande] Demande commerciale client + lignes de marchandise.
  \item[Dossier de transport] Unité transportable / facturable dérivée d'une commande confirmée.
  \item[Segment] Arrêt de chargement ou déchargement sur un dossier.
  \item[Voyage] Exécution physique (véhicule, remorque optionnelle, chauffeurs, trajets) pour un ou plusieurs dossiers.
  \item[Affectation] Lien chauffeur $\leftrightarrow$ voyage (titulaire ou renfort).
  \item[Prise de carburant] Ticket réel (litres, montant, station, statut BROUILLON / VALIDEE).
  \item[Plan d'entretien] Planning préventif lié à un engin.
  \item[Ordre de travail] Exécution atelier.
  \item[Score de santé] Indicateur saisi ; statut dérivé serveur.
  \item[Copilote IA] Assistant conversationnel via la façade Spring.
  \item[Calcul d'itinéraire IA] Estimation distance / durée via OSRM (façade Spring).
  \item[App Shell] Chrome fixe (sidebar + contenu) autour des écrans authentifiés.
  \item[File du jour] File de travail priorisée sur le tableau de bord.
\end{description}

\chapter{Décisions d'architecture (ADR)}
\label{chap:adr}

Cette annexe résume les décisions structurantes documentées dans le dépôt backend (\texttt{docs/adr/}).

\begin{description}[style=nextline,leftmargin=1.2cm,font=\normalfont\bfseries]
  \item[ADR-0001 --- Monolithe modulaire]
    Un seul déploiement Spring, plusieurs modules Modulith aux frontières \texttt{api}, plutôt qu'une myriade de microservices dès le départ.

  \item[ADR-0002 --- MFA Keycloak]
    La MFA est entièrement gérée par l'IdP ; Spring valide la présence des claims attendus dans le JWT, sans stocker d'OTP en base TMS.

  \item[ADR-0003 --- Station $\neq$ Site]
    Une station carburant n'est pas un site de chargement/déchargement : schémas et écrans séparés.

  \item[ADR-0004 --- Copilote + outils]
    Le copilote vit dans Flask ; les outils métier rappellent Spring. Base IA dédiée (pgvector) pour la mémoire agent.

  \item[ADR-0005 --- Agent planification]
    Flask propose ; Spring reste juge de la conformité et de la persistance du voyage.

  \item[ADR-0006 --- Module maintenance]
    Plans, OT, sinistres, scores et prestataires regroupés ; la flotte reste propriétaire du statut engin.
\end{description}

\noindent\textit{Références détaillées :} fichiers Markdown du dépôt \repo{logiflow-backend}/docs/adr/.


\printbibliography[heading=bibintoc,title={Bibliographie}]

\end{document}
